% Generated draft: requires mathematical and source review.
\documentclass[11pt]{article}
\usepackage[margin=25mm]{geometry}
\usepackage{fontspec}
\setmainfont{Latin Modern Roman}
\usepackage{amsmath,amssymb,mathtools,unicode-math}
\setmathfont{Latin Modern Math}
\usepackage{enumitem,needspace,xurl,hyperref}
\setlength{\parindent}{0pt}
\setlength{\parskip}{5pt}
\setlength{\emergencystretch}{3em}
\newcommand{\N}{\mathbb{N}}
\newcommand{\Z}{\mathbb{Z}}
\newcommand{\Q}{\mathbb{Q}}
\newcommand{\R}{\mathbb{R}}
\newcommand{\C}{\mathbb{C}}
\newcommand{\F}{\mathbb{F}}
\newcommand{\A}{\mathbb{A}}
\newcommand{\PP}{\mathbb{P}}
\newcommand{\E}{\mathbb{E}}
\newcommand{\eps}{\varepsilon}
\newcommand{\dd}{\,\mathrm{d}}
\begin{document}
\section*{UnsolvedMath ID: 30002742}
\textbf{Source problem code:} catalogue entry
\subsection{Definitions and mathematical statement}
Let \(G=(V,E)\) denote a graph associated with a Penrose tiling, or with another higher-dimensional quasicrystal model. The source does not specify the vertices, adjacency rule, model class, or meaning of higher-dimensionality. It also does not specify the Laplacian convention, weights, vertex measure, operator domain, or boundary conditions.

For a specified graph and a specified self-adjoint realization \(\Delta_G\) of its Laplacian on a Hilbert space \(H_G\), define
\[
\sigma(\Delta_G)=\{\lambda\in\R:\Delta_G-\lambda I\text{ does not have an everywhere-defined bounded inverse on }H_G\}.
\]
The pure point, absolutely continuous, and singular continuous spectral subspaces are the subspaces determined by the corresponding types of spectral measures of \(\Delta_G\). An eigenfunction at energy \(\lambda\) is a nonzero \(f\in\operatorname{Dom}(\Delta_G)\) satisfying \(\Delta_G f=\lambda f\); it has finite support if \(\{v\in V:f(v)\ne0\}\) is finite.

\textbf{Problem, conditional on specifying the model and operator.} Determine the spectral properties of \(\Delta_G\) for a graph defined by a Penrose tiling, or by another higher-dimensional quasicrystal model. This includes characterizing its spectrum and spectral type. The supplied dated assessment additionally identifies spectral gaps and transport as incompletely understood topics, but gives no precise gap question, transport observable, initial state, or limiting regime.

This is a research program rather than a uniquely specified conjecture. The source does not determine whether the requested conclusions concern one chosen tiling, every tiling in a specified family, or an almost-everywhere statement for a specified probability measure. No such quantifier, graph construction, or operator convention is selected here. Questions about an integrated density of states likewise require a definition, normalization, and limiting procedure that the supplied record does not provide.
\subsection{Short English statement}
Determine the spectrum and spectral type of a specified Laplacian on a graph obtained from a Penrose tiling or another higher-dimensional quasicrystal model. The source leaves the graph construction, operator, quantifiers, and precise spectral targets unresolved, so this remains a conditional statement requiring mathematical review.
\subsection{Sources}
\begin{itemize}
\item Local registry: \texttt{data/unsolvedmath/problems.json}, permanent UnsolvedMath ID \(30002742\), source problem code \texttt{catalogue entry}, title \emph{Laplacian Spectra of Penrose Tiling Graphs}. The record belongs to the Oberwolfach Reports Open Problems collection; no individual proposer is supplied.
\item Supplied source citation: \emph{Mini-Workshop: Dynamical versus Diffraction Spectra in the Theory of Quasicrystals} (2015), Oberwolfach Reports, pp.\ 2987--3013, DOI \url{https://doi.org/10.4171/owr/2014/53}. The background describes Oberwolfach Reports as documenting mathematical workshops and their open-problem discussions. This is also the literature-source entry labelled ``original OWR report.'' The registry separately gives proposed year 2014.
\item Supplied literature source: Damanik et al., \emph{Laplacians on Penrose tilings}, \url{https://doi.org/10.1080/10586458.2023.2206589}.
\item Supplied literature source labelled ``geometric features of spectra'': \url{https://doi.org/10.1007/s10208-022-09598-0}. No fuller bibliographic attribution is supplied.
\item Literature assessment dated exactly \texttt{2026-08-21}: the supplied curation labels the problem \texttt{partially\_solved} and reports results on locally supported eigenfunctions and discontinuities in the integrated density of states for Penrose-related Laplacians, while describing overall spectral type, gaps, and transport as incompletely understood. This is supplied dated curation derived from the OWR record, not an independently verified assessment. The supplied background expressly requests independent verification of time-sensitive status and cited claims.
\end{itemize}
\end{document}
